\section{Research Context}
\label{sec:research-context}

This thesis is situated within the broader research initiatives of the University of Canterbury Computer Vision Research Group (UCVision), specifically contributing to the MBIE funded project, `Predicting the unseen: a new method for accurate yield estimation in viticulture/horticulture'. This project utilises autonomous field robots (\cref{fig:robot-apple-orchard}) to capture multi-view imagery of vineyards and orchards, which are subsequently processed \textit{offline} to generate high-fidelity 3D reconstructions using 3D Gaussian Splatting (3DGS). These models serve as the ``digital twin'' required for complex agricultural tasks such as precision yield estimation and autonomous intervention.

The primary objective of this project is the automation of vineyard pruning. In this workflow, pruning decisions and robotic arm path planning are computed within the high-fidelity 3DGS environment. Unlike conventional pruning systems that rely on lower-resolution, real-time sensor data for immediate modelling and execution, the UCVision workflow decouples data collection from intervention. While this allows for superior structural analysis of the vine, it introduces a fundamental spatial synchronisation problem: the alignment of the robot's current pose to the offline model's coordinate frame.

For the robot to execute a pre-computed pruning plan, it must map its real-time physical state to the coordinate system of the offline model. This necessitates a robust SLAM framework capable of solving three distinct sub-problems:
\begin{enumerate}
    \item \textbf{Global Relocalisation:} Navigating the platform back to the specific plant or row identified in the modelling phase.
    \item \textbf{Strategic Positioning:} Aligning the mobile base within the physical workspace constraints of the robotic arm to ensure all required cutpoints are reachable.
    \item \textbf{Sim-to-Real Registration:} Achieving centimetre-level alignment between the robot's onboard perception and the 3D model, ensuring that the path planner's simulated coordinates correspond precisely to the physical vine structures.
\end{enumerate}

To the author's knowledge, no existing vineyard robotics system addresses this full registration chain. Current pruning platforms assume a tightly coupled perception--execution loop in which the model is built and consumed in real time; they do not face the problem of aligning a robot to a pre-built offline reconstruction. Similarly, existing agricultural SLAM systems target navigation-grade accuracy and do not provide the centimetre-level pose refinement required for manipulation against an offline model.

This thesis addresses this capability gap by developing a SLAM and relocalisation framework that enables the registration chain described above. By synchronising the robot’s perceived position with both the physical environment and the pre-existing 3D model, this research enables the transition from static digital twins to active, autonomous robotic intervention in the field.